

\cvsection{Academic positions and education}

%\textit{\footnotesize My academic career spans positions across Mathematics and Computer Science, with a unifying focus on dualities in logic: from commutator theory in universal algebra, to domain theory in denotational semantics, and to non-classical logics.}



%-------------------------------------------------------------------------------
\cvsubsection{\stress{Academic positions}}
%-------------------------------------------------------------------------------

\begin{cventries}
	
	%---------------------------------------------------------
	
	
	\cventry
	{FSR Incoming Postdoctoral Fellowship} % Degree
	{Université catholique de Louvain (Belgium), Research Institute in Mathematics and Physics.} % Institution
	{2025.09.01 -- present} % Location
	{} % Date(s)
	{\begin{cvitems} % Description(s) bullet points
%			\item %Fellowship: FSR Incoming Post-Doc Fellowship. 
%			Awarded  by the Université catholique de Louvain.
			\item Project: \emph{CatCoD: Categorical Commutators and Dualities}.
			\item Supervisor: Marino Gran.
	\end{cvitems}}
	%---------------------------------------------------------
	
	\bigskip\smallskip
	
	\cventry
	{Marie Skłodowska--Curie Postdoctoral Fellowship (MSCA PF)} % Degree
	{University of Birmingham (UK), School of Computer Science.} % Institution
	{2023.09.01 -- 2025.08.31} % Location
	{} % Date(s)
	{\begin{cvitems} % Description(s) bullet points
			\item Awarded after EU MSCA evaluation; funded by UK Research and Innovation under the Horizon Europe funding guarantee.
%			Marie Sk{\l}odowska Curie Postdoctoral Fellowship.  Awarded by the European Union under the Horizon 2020 Programme, and delivered by UK Research and Innovation under the Horizon Europe funding guarantee (grant number EP/Y015029/1).
%			\item Host institution: School of Computer Science, University of Birmingham, within the Theory of Computation group.
			\item Project: \emph{DCPOS: Dualities via ComPactness, Openness and their Symmetry}.
			\item Supervisor: Achim Jung.
	\end{cvitems}}
	%---------------------------------------------------------
	
	\bigskip\smallskip
  
  \cventry
  {PostDoc} % Degree
  {University of Salerno (Italy), Department of Mathematics.} % Institution
  {2020.12.01 -- 2023.08.31} % Location
  {} % Date(s)
  {\begin{cvitems} % Description(s) bullet points
  		\item {Project: \emph{{\L}ukasiewicz logic, MV-algebras and duality}.}
  		\item Supervisor: Luca Spada.
  \end{cvitems}}
  %---------------------------------------------------------
  
\end{cventries}


\bigskip\bigskip


\cvsubsection{\stress{Education}}

\begin{cventries}
      
  \cventry
    {PhD {\enskip\cdotp\enskip} Mathematics} % Degree
    {University of Milan (Italy), Department of Mathematics.} % Institution
    {2017.09.25 -- 2021.04.09} % Location
    {} % Date(s)
    {
      \begin{cvitems} % Description(s) bullet points
%      	\item Academic title awarded: ``Dottore di Ricerca in Scienze Matematiche'' (Degree of Doctor of Philosophy in Mathematical Sciences).
        \item {Thesis: \emph{On the axiomatisability of the dual of compact ordered spaces}.}
        \item {Supervisor: Vincenzo Marra.} % (Department of Mathematics, University of Milan, Italy).}
%        \item {Defense: 9 April 2021.}
        \item {Grade: Excellent with honours.}
%        \item {PhD defense committee: 
%        	%
%        	\begin{enumerate}
%        		\item Giuseppe Rosolini (Chairman -- University of Genoa),
%        		\item Samuel J.~van~Gool (Secretary -- IRIF, Université Paris-Cité),
%        		\item Guram Bezhanisvili (New Mexico State University).
%        	\end{enumerate}
%        	%
%        	}
%        \item Three years, fully funded by the University of Milan.
        \item Abstract published in the Bulletin of Symbolic Logic. DOI: \href{https://doi.org/10.1017/bsl.2021.54}{10.1017/bsl.2021.54}.
%        \item {Content: We prove that the category of Nachbin's compact ordered spaces and order-preserving continuous maps between them is dually equivalent to a variety of algebras, with operations of at most countable arity.
%        This solves an open problem in a paper by D.~Hofmann, R.~Neves and P.~Nora (2018).
%       	We exhibit a \emph{finite} equational axiomatisation of the dual of the category of compact ordered spaces, meaning that our description uses only finitely many function symbols and finitely many equational axioms.
%       	In preparation for the latter result, we generalise a celebrated theorem by D.~Mundici (1986) from a categorical equivalence for unital Abelian lattice-ordered groups to a categorical equivalence for unital commutative distributive lattice-ordered \emph{monoids}.}
%       \item Cited by:
%       %
%       \begin{enumerate}
%       	\item 
%       	S.~v.~Gool, J.~Marquès, On duality and model theory for polyadic spaces. Preprint at arXiv:2210.01018, 2022.
%       	\item
%       	M.\ Abbadini, P.\ Jipsen, T.\ Kroupa, S.\ Vannucci.
%       	A finite axiomatization of positive MV-algebras.
%       	{\em Algebra Universalis}, 83:28, 2022.
%       	\item
%       	M.\ Abbadini.
%       	Equivalence à la Mundici for commutative lattice-ordered monoids.
%       	{\em {Algebra Universalis}}, 82:45, 2021.
%       	DOI: \href{https://doi.org/10.1007/s00012-021-00736-3}{10.1007/s00012-021-00736-3}.
%       \end{enumerate}
%       %
      \end{cvitems}
    }
  
%  \bigskip\smallskip

  \cventry
    {MSc{\enskip\cdotp\enskip}Mathematics} % Degree
    {University of Milan (Italy), Department of Mathematics.} % Institution
    {2015.11.01 -- 2017.07.19} % Location
    {} % Date(s)
    {%
%      \begin{cvitems} % Description(s) bullet points
%      	\item Academic title awarded: ``Dottore Magistrale in Matematica'' (Master Degree in Mathematics).
%        \item {Thesis: \emph{The algebraic structure of probability measures and expectations}.}
%%        \item {Topics: MV-algebras, lattice-ordered groups and their applications to measure theory.}
%        \item {Supervisor: Vincenzo Marra.}
%%        \item {Date of graduation: 19 July 2017.}
%        \item {Grade: 110/110 cum laude.}
%        \item {Average grade: 29.63/30.}
%	     	\item {Erasmus+: Semester abroad at the University of Bonn, Germany (Spring 2015/16).}
%%	     \vspace{0.5mm}
%	      \item {Erasmus+: Semester abroad at the University of Stuttgart, Germany (Fall 2016/17).}
%	    \end{cvitems}
	  }
	  
	  \vspace{-1.2em}
  
%  \bigskip\smallskip

  \cventry
    {BSc{\enskip\cdotp\enskip}Mathematics} % Degree
    {University of Milan (Italy), Department of Mathematics.} % Institution
    {2012.07.16 -- 2015.10.21} % Location
    {} % Date(s)
    {%
%       \begin{cvitems} % Description(s) bullet points
%       	\item Academic title awarded: ``Dottore in Matematica'' (Bachelor Degree in Mathematics).
%%        \item {Thesis: \emph{Peano-Lawvere axioms, decidable sets, and Goodstein's Theorem}.}
%%        \item {Advisor: Prof.\ Emanuele Pacifici (Department of Mathematics, Univesity of Milan, Italy).}
%%        \item {Final grade: 106/110.}
%%        \item {Date of graduation: 21 October 2015.}
%%         \item {I have written the manuscript `\emph{Dealing with Permutation Puzzles independently}', on applications of group theory to the Rubik's cube and similar puzzles, under the supervision of Prof.\ Emanuele Pacifici.}
%       \end{cvitems}
    }
%---------------------------------------------------------
\end{cventries}

	  \vspace{-1.2em}
